\documentclass[11pt]{article}

\usepackage[a4paper,margin=0.9in]{geometry}
\usepackage{amsmath,amssymb}
\usepackage[colorlinks=true,allcolors=blue]{hyperref}
\usepackage{parskip}
\setlength{\parskip}{0.5em}

\begin{document}

\thispagestyle{empty}

\begin{flushright}
July 2026
\end{flushright}

\vspace{-1em}

\noindent
The Editor-in-Chief, \textit{IEEE Internet of Things Journal}

\vspace{0.6em}

\noindent
Dear Editor,

We are pleased to submit our manuscript entitled ``Congestion-Aware Joint
Charging-Station Placement, Capacity, and Swarm-UAV Trajectory Design for
Age-Optimal Data Collection'' for consideration for publication in the
\textit{IEEE Internet of Things Journal}.

Unmanned aerial vehicles increasingly collect time-sensitive Internet of Things
data, where the quality metric is the Age of Information. Practical deployments
use a swarm of energy-limited UAVs that recharge at a few shared charging
stations, which then become a contended resource whose waiting time inflates the
Age of Information. Jointly deciding where to place the stations, how many
charging ports to install, how the swarm should fly, and which UAV recharges
where is a coupled problem that, to our knowledge, has not been addressed as a
whole.

Our contributions are: (i) the first joint formulation of charging-station
placement, port capacity, swarm trajectory, and assignment for peak-Age-of-Information
data collection with a contended, multi-server charging resource; (ii) a
finite-source (machine-repair) queue model, validated against a discrete-event
simulation, which shows that the commonly assumed open $M/M/c$ queue overpredicts
the waiting time by roughly six times and would distort the design; (iii) two
proved structural results, a congestion threshold that makes the Age of
Information non-monotone in the swarm size with a provisioning law, and a
placement inversion from coverage-driven to traffic-driven; and (iv) a
block-coordinate solver with a GPU close-enough traveling-salesman trajectory
optimizer, evaluated against several baselines with large-swarm scaling,
robustness, and optimality-gap studies.

The topic aligns with the scope of the journal and will interest its readership
working on UAV-assisted data collection, Age of Information, and energy-aware
Internet of Things systems.

We confirm that the manuscript is original, has not been published previously,
and is not under consideration elsewhere; all authors approved the submission and
declare no conflict of interest. To support reproducibility, the source code and
data are available at \url{https://github.com/haodpsut/swarm-uav-aoi-congestion}.

Thank you for considering our manuscript.

\vspace{0.4em}

\noindent
Sincerely,\\
Truong Duy Dinh (corresponding author), on behalf of\\
Phuc Hao Do, Tran Duc Le, Truong Duy Dinh, and Van Dai Pham\\
\href{mailto:duydt@ptit.edu.vn}{duydt@ptit.edu.vn}

\end{document}
